% Integración por delta de la adenda de realizaciones físicas.
% Residencia principal: transferencia física de las tres álgebras y lector temporal.

\chapter{Régimes quadratiques et lecture temporelle}
\label{chap:regimenes-cuadraticos-tiempo}
\index{algèbre quadratique!trois régimes}
\index{cristal temporel!réalisation finie HMT}
\index{temps!lecteur TRIT}

\section{Algèbres quadratiques du TRIT et signatures de courbure}
\label{sec:regimenes-cuadraticos-recibidos}

La construction et la démonstration primaires des trois régimes résident dans
le \cref{chap:trit-algebras}. Pour
\(\tau\in\TRIT\), avec une image équilibrée
\(\bar\tau\in\{-1,0,+1\}\), on a
\[
 \mathbb A_\tau=\RR[X]/(X^2+\bar\tau),
\]
et, si \(J\) est la classe de \(X\),
\[
\begin{array}{ccl}
 \tau=+1&:&J^2=-I,\quad\text{régime elliptique},\\
 \tau=0&:&J^2=0,\quad\text{régime parabolique},\\
 \tau=-1&:&J^2=+I,\quad\text{régime hyperbolique}.
\end{array}
\]
Pour les trois valeurs de \(\tau\), le quotient
\(\mathbb A_\tau=\mathbb R[X]/(X^2+\bar\tau)\) satisfait exactement les trois
identités quadratiques précédentes. En particulier, la racine interne de
\(-1\) se réalise comme opérateur dans le feuillet elliptique ; HMT n’a besoin ni de nier ni de
supprimer l’extension complexe. L’algèbre parabolique contient un nilpotent,
et l’algèbre hyperbolique se décompose au moyen des deux idempotents
\((I\pm J)/2\).

Ces identités classent des algèbres et des transports. Elles ne sélectionnent pas à elles
seules une métrique spatio-temporelle et ne prouvent pas que l’univers physique occupe l’un
des trois feuillets.

\section{Famille métrique paramétrée du lecteur temporel}
\label{sec:ansatz-metrico-tres-curvaturas}

Écrivons
\[
 \mathbb A_s:=\RR[X]/(X^2+s),
 \qquad s\in\{+1,0,-1\},
\]
pour la même famille avec le signe réel comme indice.
Pour préciser l'interface possible, fixons une échelle
\(\rho_{\rm curv}>0\), de dimension \(L^{-1}\), et
\(\theta\in\RR/2\pi\ZZ\). Pour \(s\in\{+1,0,-1\}\), soit
\[
 g_{s,\rho}
 =\dd r^2+S_{s,\rho}(r)^2\,\dd\theta^2,
\]
\[
 S_{s,\rho}(r)=
 \begin{cases}
  \rho_{\rm curv}^{-1}\sin(\rho_{\rm curv}r),&s=+1,\\
  r,&s=0,\\
  \rho_{\rm curv}^{-1}\sinh(\rho_{\rm curv}r),&s=-1.
 \end{cases}
\]
Le domaine radial est \(r\geq0\) pour \(s=0,-1\) ; pour \(s=+1\), la carte
polaire couvre \(0\leq r\leq\pi/\rho_{\rm curv}\), avec la régularité polaire
usuelle aux extrémités. À l'intérieur de chaque carte,
\[
 K_{\rm G}
 =-\frac{S_{s,\rho}''}{S_{s,\rho}}
 =s\,\rho_{\rm curv}^{\,2}.
\]
La dernière égalité est un calcul exact du modèle une fois la
réalisation fixée. La flèche
\[
 \mathbb A_s\dashrightarrow(M_s,g_{s,\rho})
\]
est cependant un \emph{ansatz} d'\textsc{interface physique} ; ce n'est pas une
conséquence canonique de la relation \(J^2=-sI\).

Si \(\ell_0\in\mathcal L_L^+\) est une longueur et que l'on déclare
\[
 \rho_{\rm curv}:=\frac{\Delta_4}{\ell_0},
\]
alors
\[
 K_{\rm G}
 =s\,\frac{\Delta_4^2}{\ell_0^2}.
\]
Cette spécialisation est
\textsc{démontrée avec une structure de départ explicite} : la conclusion est
obtenue une fois \(s,\Delta_4,\ell_0\) déclarés, mais elle ne sélectionne pas
\(\ell_0\), ne
déduit pas la courbure cosmologique observée et n'identifie pas les rapports
\(\mathcal G_{\rm tor}/\mathcal I_{\rm tor}\) ou
\(\mathcal G_{\rm av}/\mathcal I_{\rm av}\) à \(K_{\rm G}\).

\section{Lecteur temporel du TRIT sur ({-1,0,+1}) : ouverture, seuil et retour avec mémoire}
\label{sec:lector-temporal-trit}

Une fois choisie une paramétrisation ordonnée des transitions, le lecteur
temporel actif est
\[
\begin{aligned}
 +1&\longmapsto
 \text{retour, mémoire retenue et passé},\\
 0&\longmapsto
 \text{seuil d'évaluation et présent},\\
 -1&\longmapsto
 \text{ouverture bidirectionnelle et futur}.
\end{aligned}
\]
Le lecteur est \textsc{démontré avec une structure de départ explicite} : ses
conclusions s’obtiennent une fois déclarés la signature TRIT, l’ordre de la
trajectoire et la mémoire transportée. Il ne transforme pas l’ensemble
\(\{-1,0,+1\}\) en une équation du mouvement. Le futur désigne l’espace
des prolongements compatibles ; la survie peut sélectionner une branche
ou une multisection sans effacer la généalogie des autres.

La fréquence
\[
 \omega_0=\frac{2\pi}{108\,t_0}
\]
type l’horloge du retour complet au moyen de l’entropie de décision
\(I_{\mathrm{TRIT}}=\log 3\) et du transducteur de Boltzmann.
Elle ne démontre pas, isolément, une phase sous-harmonique stable. Elle permet en revanche de construire le
Hamiltonien fini de référence qui suit ; la sélection d’un Hamiltonien
physique robuste requiert des conditions supplémentaires.

\section{Dynamique unitaire sur 108 états : Hamiltonien, propagateur et période}
\label{sec:hamiltoniano-reloj-108}

Soit
\[
 \mathcal H_{\rm clk}:=\mathbb C^{108},
 \qquad
 \{|j\rangle:j\in\mathbb Z/108\mathbb Z\}
\]
muni de son produit hermitien canonique. Fixons
\(\zeta:=\exp(2\pi\mathrm i/108)\), le décalage unitaire
\[
 U_{\rm clk}|j\rangle=|j+1\rangle
\]
et la base de Fourier
\[
 |\widetilde k\rangle
 :=\frac1{\sqrt{108}}\sum_{j=0}^{107}\zeta^{jk}|j\rangle,
 \qquad 0\leq k\leq107.
\]
Avec \(t_0\in\mathcal L_T^+\) et
\(\hbar_{\rm int}\in\mathcal L_S^+\), nous définissons
\[
 \boxed{
 H_{\rm clk}
 :=\hbar_{\rm int}\omega_0
 \sum_{k=0}^{107}k\,
 |\widetilde k\rangle\langle\widetilde k|,
 \qquad
 \omega_0=\frac{2\pi}{108\,t_0}.}
 \tag{CLK1}\label{eq:hamiltoniano-reloj-108}
\]
L’opérateur a pour type énergétique
\([H_{\rm clk}]=\mathcal L_S\mathcal L_T^{-1}=\mathcal L_E\).

\begin{theorem}[Réalisation unitaire exacte du calendrier \(108\)]
\label{thm:reloj-unitario-108}
L’opérateur \(H_{\rm clk}\) est autoadjoint et son propagateur pendant un pas
satisfait
\[
 U_{\rm step}
 :=\exp\!\left(-\frac{\mathrm i}{\hbar_{\rm int}}
 H_{\rm clk}t_0\right)
 =U_{\rm clk},
 \qquad
 U_{\rm step}^{108}=I.
 \tag{CLK2}\label{eq:propagador-reloj-108}
\]
Si
\[
 Z_{\rm clk}|j\rangle=\zeta^j|j\rangle,
\]
alors, pour tout état de base \(|j_0\rangle\),
\[
 \langle Z_{\rm clk}\rangle_{n}
 :=\langle j_0|U_{\rm step}^{-n}Z_{\rm clk}
 U_{\rm step}^n|j_0\rangle
 =\zeta^{j_0+n}.
 \tag{CLK3}\label{eq:respuesta-reloj-108}
\]
La suite a pour période primitive \(108\).
\end{theorem}

\begin{proof}
L’expression \eqref{eq:hamiltoniano-reloj-108} est une décomposition
spectrale à valeurs réelles, donc \(H_{\rm clk}=H_{\rm clk}^\dagger\).
En outre,
\[
 U_{\rm clk}|\widetilde k\rangle
 =\zeta^{-k}|\widetilde k\rangle,
\]
tandis que l’exponentielle de
\(\hbar_{\rm int}\omega_0 k\) pendant \(t_0\) produit
\(\exp(-2\pi\mathrm ik/108)=\zeta^{-k}\). Les deux opérateurs coïncident sur
une base orthonormale. Les identités restantes découlent de
\(U_{\rm step}^n|j_0\rangle=|j_0+n\rangle\) et de la primitivité de
\(\zeta\).
\end{proof}

La formulation variationnelle correspondante, pour une courbe
\(\psi\in C^1([t_a,t_b],\mathcal H_{\rm clk})\), non contrainte pendant la
variation, est
\[
 \boxed{
 L_{\rm clk}(\psi,\dot\psi)
 =\frac{\mathrm i\hbar_{\rm int}}2
 \bigl(
 \langle\psi|\dot\psi\rangle
 -\langle\dot\psi|\psi\rangle
 \bigr)
 -\langle\psi|H_{\rm clk}|\psi\rangle.}
 \tag{CLK4}\label{eq:lagrangiano-reloj-108}
\]
La fonctionnelle d’action est
\[
 \boxed{
 S_{\rm clk}[\psi]
 :=\int_{t_a}^{t_b}
 L_{\rm clk}(\psi(t),\dot\psi(t))\,\mathrm dt,}
 \tag{CLK5}\label{eq:accion-reloj-108}
\]
avec \(\delta\psi(t_a)=\delta\psi(t_b)=0\).
\index{Schrödinger!équation de l’horloge finie}
Faire varier \(\psi\) et \(\bar\psi\) indépendamment produit l’équation de
Schrödinger de l’horloge finie
\[
 \boxed{\mathrm i\hbar_{\rm int}\dot\psi=H_{\rm clk}\psi}
 \tag{CLK6}\label{eq:schrodinger-reloj-108}
\]
ainsi que son équation adjointe \cite{SchrodingerNobel}. Son propagateur échantillonné
en \(t=nt_0\) est \(\psi_n=U_{\rm step}^{n}\psi_0\). Comme \(H_{\rm clk}\) est
autoadjoint, la norme est conservée ; une donnée initiale normalisée reste
normalisée. Le résultat est exact par rapport aux
\(H_{\rm clk},t_0,\hbar_{\rm int}\) déclarés et ne sélectionne pas de Hamiltonien
physique ni n’introduit de fréquence expérimentale.

Le \cref{thm:reloj-unitario-108} démontre la dynamique unitaire du calendrier
relativement à la complétion complexe, à \(t_0\) et à \(\hbar_{\rm int}\)
déclarés. Les cinq conditions dynamiques de la section suivante reçoivent
une réalisation constructive dans \cref{sec:hmtf2-cristal-temporal} : on y
fixe l'entraînement, l'observable d'ordre, la sous-harmonie primitive de période
\(108T_{\rm d}\) et la stabilité spectrale finie. La réalisation
macroscopique matérielle constitue une reconnaissance externe distincte et non une
lacune du résultat fini.

\section{Conditions dynamiques de réalisation d'un cristal temporel : hamiltonien, réponse sous-harmonique et stabilité}
\label{sec:programa-cristal-tiempo}

Une réalisation entraînée doit fournir, au minimum :
\begin{enumerate}
 \item un espace d'états physique et un hamiltonien autoadjoint
 \(H(t)\), ou un protocole entraîné équivalent, avec
 \(H(t+T_{\rm d})=H(t)\) ;
 \item le propagateur de Floquet
 \[
  U_F=\mathcal T\exp\!\left(
  -\frac{\mathrm i}{\hbar_{\rm int}}
  \int_0^{T_{\rm d}}H(t)\,\dd t\right);
 \]
 \item un observable d'ordre \(\mathcal O\) et une famille d'états pour
 lesquels il existe une réponse sous-harmonique primitive d'ordre \(m>1\) :
 \[
  \langle\mathcal O((n+m)T_{\rm d})\rangle
  =\langle\mathcal O(nT_{\rm d})\rangle
 \]
 dans le régime asymptotique, sans la même périodicité pour un diviseur propre
 de \(m\) ;
 \item la stabilité de \(m\), de la phase et de l'amplitude face à une classe
 ouverte de perturbations admissibles ;
 \item une persistance qui ne soit pas un transitoire de taille finie, avec un
 mécanisme déclaré de localisation, de préthermalisation ou de protection
 dynamique.
\end{enumerate}
Le calendrier \(27\to54\to108\) fournit la hiérarchie interne des périodes.
La réalisation finie de \cref{sec:hmtf2-cristal-temporal} sélectionne l'ordre
sous-harmonique \(108\) dans son domaine déclaré, sans remplacer les cinq
constructions précédentes ni imposer à elle seule une phase matérielle
macroscopique.

\begin{criterion}[Cristal temporel HMT--MD]
\label{crit:cristal-tiempo-hmt-md}
L'identification est réfutée s'il n'existe pas de hamiltonien ou de protocole
entraîné bien typé ; si la sous-harmonie supposée disparaît sous des
perturbations arbitrairement petites compatibles avec le modèle ; si la
période observée n'est que la période imposée par l'impulsion ; si
l'oscillation est un transitoire fini ; ou si le lecteur temporel assigne des histoires
équivalentes à des phases physiques incompatibles.
\end{criterion}

\paragraph{Domaine mathématique et physique du lecteur temporel.}
Les trois algèbres et leurs exponentielles sont démontrées dans le domaine algébrique ;
le lecteur temporel se déduit de la signature, de l'ordre et de la mémoire initiaux. Le
cristal temporel fini HMT satisfait le
\cref{crit:cristal-tiempo-hmt-md} au moyen de la construction de
\cref{sec:hmtf2-cristal-temporal}. L'extension à une phase matérielle
macroscopique exige, en outre, de spécifier le support physique et relève de la
confrontation expérimentale ultérieure.
